\documentclass[14pt,a4paper]{extarticle}
\usepackage[T2A]{fontenc}
\usepackage[utf8]{inputenc}
\usepackage[english,russian]{babel}
\usepackage{lmodern}
\usepackage{geometry}
\usepackage{array}
\usepackage{tabularx}
\usepackage{longtable}
\usepackage{ltablex}
\usepackage{booktabs}
\usepackage{hyperref}
\hypersetup{hidelinks}
\usepackage{indentfirst}
\usepackage{microtype}
\usepackage{graphicx}
\usepackage{setspace}
\usepackage{listings}
\usepackage{graphicx}
\usepackage{caption}
\usepackage{float}
\usepackage{tocloft}
\usepackage{titlesec}
\usepackage{newtxtext,newtxmath}

\titleformat{\section}        % форматируем \section
  {\normalfont\bfseries}      % стиль: обычный + жирный
  {\thesection.}              % номер секции с точкой
  {2mm}                       % расстояние между номером и текстом
  {}                           % перед заголовком можно добавить что-то
\titlespacing*{\section}{1.5cm}{*3}{*1}  % {отступ слева}{отступ сверху}{отступ снизу}

\titleformat{\subsection}        % форматируем \subsection
  {\normalfont\bfseries}      % стиль: обычный + жирный
  {\thesubsection.}              % номер секции с точкой
  {2mm}                       % расстояние между номером и текстом
  {}                           % перед заголовком можно добавить что-то
\titlespacing*{\subsection}{1.5cm}{*3}{*1}  % {отступ слева}{отступ сверху}{отступ снизу}

\captionsetup{justification=centering}
\captionsetup[table]{labelformat=empty,justification=raggedleft,singlelinecheck=false}
\renewcommand{\ttdefault}{cmtt}
\keepXColumns
\setlength{\parindent}{3em}
\newgeometry{left=2.5cm,right=2.5cm,top=3cm,bottom=2.5cm}
\graphicspath{{../}{./}}

\renewcommand{\contentsname}{\hspace*{4em}{\smallСодержание}}
\setlength{\cftbeforesecskip}{0.5em}
\renewcommand{\cftsecleader}{\cftdotfill{\cftdotsep}}
\renewcommand{\cftsecaftersnum}{.\hspace{2mm}}

\title{Отчёт по "Корпоративные информационные системы"}
\author{Выполнил студент группы М-РИТ-25-1, Станиславчук Сергей}

\begin{document}

\setstretch{1.5}
\noindent
\begin{minipage}[c]{1.2cm}
    \includegraphics[width=1.1cm]{~/programmer/++Programmer/5_1/LSTU.jpg}
\end{minipage}%
\begin{minipage}[c]{0.92\linewidth}
    \centering
    {\small
    МИНИСТЕРСТВО НАУКИ И ВЫСШЕГО ОБРАЗОВАНИЯ РФ \\[1em]
    ФЕДЕРАЛЬНОЕ ГОСУДАРСТВЕННОЕ БЮДЖЕТНОЕ \\
    ОБРАЗОВАТЕЛЬНОЕ УЧРЕЖДЕНИЕ ВЫСШЕГО ОБРАЗОВАНИЯ \\
    «ЛИПЕЦКИЙ ГОСУДАРСТВЕННЫЙ ТЕХНИЧЕСКИЙ УНИВЕРСИТЕТ»
    }
\end{minipage}%

\vspace{3em}

\noindent
\hspace{1em}%
\begin{minipage}[t]{0.8\textwidth}
\textbf{Институт} \hspace{2em} компьютерных наук \\
\textbf{Кафедра} \hspace{2em} автоматизированных систем управления
\end{minipage}


\vspace{6em}

\begin{center}
\textbf{Практическая работа №2} \\
по дисциплине \\ <<Корпоративные информационные системы>>
\end{center}

\vspace{6em}

\noindent
\begin{tabular*}{\textwidth}{@{\extracolsep{\fill}} l c r }
    Студент М-РИТ-25-1 & \makebox[4cm]{\hrulefill} & Станиславчук С.М. \\[2em]
Руководитель & \\
К.т.н., доцент & \makebox[4cm]{\hrulefill} & Галкин А.В. \\
\end{tabular*}

\vfill
\begin{center}
    Липецк, 2026
\end{center}

\thispagestyle{empty}

\newpage
\tableofcontents
\newpage

\section{Задание}

Цель работы --- разработать в среде «1С:Предприятие 8.3» конфигурацию для
автоматизации учёта успеваемости студентов высшего учебного заведения.
Объектом автоматизации является учебный процесс: хранение сведений о студентах,
преподавателях, учебных группах, дисциплинах и семестрах, регистрация оценок,
их корректировка, а также расчёт показателей успеваемости.

Работа выполняется на основе ранее выбранной автоматизированной информационной
системы «Система управления учебным процессом в высшем учебном заведении»
(личный кабинет). В системе предусмотрены три роли пользователей: студент,
преподаватель и администратор. Администратор ведёт нормативно-справочную
информацию и назначает преподавателей, преподаватель выставляет и корректирует
оценки, студент просматривает свои результаты.

\textbf{Этап 1. Проектирование и разработка структуры конфигурации.}
Требуется создать «каркас» базы данных в конфигураторе, определить состав
объектов автоматизации, их реквизиты, типы данных и иерархию:

\begin{itemize}
    \item разработать не менее пяти справочников и описать их стандартные и
    дополнительные реквизиты;
    \item разработать не менее трёх документов, определить реквизиты шапки,
    табличных частей и показатели документов;
    \item разработать не менее двух регистров накопления, определить их вид,
    измерения, ресурсы и источники формирования движений;
    \item описать иерархическую структуру конфигурации.
\end{itemize}

\textbf{Этап 2. Реализация функционала автоматизации и программирование.}
Требуется настроить логику работы информационной системы, автоматизировать
расчёты и исключить ошибки пользователя с помощью программного кода:

\begin{itemize}
    \item запрограммировать экранные формы и логику элементов, реализовать
    автоматический расчёт значений и автоподстановку данных;
    \item разработать механизмы проведения документов и формирования движений
    по регистрам накопления;
    \item реализовать контроль и бизнес-логику, защищающие от некорректного
    ввода и проведения документов.
\end{itemize}

\section{Ход работы}

\subsection{Этап 1. Проектирование структуры конфигурации}

На первом этапе в конфигураторе создана иерархия объектов, отражающая
предметную область учёта успеваемости. В состав конфигурации вошли шесть
справочников, три документа, два регистра накопления и одно перечисление.
Количество объектов превышает минимальные требования задания: справочников~---
6 (требуется не менее 5), документов~--- 3 (не менее 3), регистров
накопления~--- 2 (не менее 2).

\noindent\textbf{Справочники.}
Для хранения условно-постоянной информации об участниках учебного процесса
и учебных объектах создано шесть справочников. Их состав и ответственные
за заполнение приведены в таблице~\ref{tab:refs}.

\begin{table}[H]
    \centering
    \caption{Таблица \thetable\ --- Состав справочников конфигурации}
    \label{tab:refs}
    \footnotesize
    \begin{tabularx}{\textwidth}{@{}l >{\raggedright\arraybackslash}X l@{}}
        \toprule
        Справочник & Назначение & Ответственный \\
        \midrule
        Студенты & Сведения о студентах: ФИО, номер зачётной книжки, дата
        рождения, учебная группа & Администратор \\
        Преподаватели & Сведения о преподавателях: ФИО, должность,
        квалификация & Администратор \\
        УчебныеГруппы & Учебные группы: название, курс, факультет &
        Администратор \\
        Дисциплины & Учебные дисциплины: название, код, количество часов &
        Администратор \\
        Семестры & Учебные периоды: учебный год, номер семестра, даты начала и
        окончания & Администратор \\
        ПользователиИС & Учётные записи пользователей: логин, роль, связь со
        студентом или преподавателем & Администратор \\
        \bottomrule
    \end{tabularx}
    \addtocounter{table}{-1}
\end{table}

Для всех справочников используются стандартные реквизиты \texttt{Код} и
\texttt{Наименование}; ниже приведены дополнительные реквизиты и уточнено
назначение наименования.

\noindent\textbf{Справочник <<Студенты>>.}
Предназначен для хранения сведений о студентах. Реквизиты:

\begin{itemize}
    \item \texttt{Наименование} --- ФИО студента (строка);
    \item \texttt{НомерЗачетнойКнижки} --- номер зачётной книжки (строка);
    \item \texttt{ДатаРождения} --- дата рождения (дата);
    \item \texttt{УчебнаяГруппа} --- ссылка на справочник <<УчебныеГруппы>>.
\end{itemize}

\noindent\textbf{Справочник <<Преподаватели>>.}
Предназначен для хранения сведений о преподавателях. Реквизиты:

\begin{itemize}
    \item \texttt{Наименование} --- ФИО преподавателя (строка);
    \item \texttt{Должность} --- должность (строка);
    \item \texttt{Квалификация} --- квалификация или учёная степень (строка).
\end{itemize}

\noindent\textbf{Справочник <<УчебныеГруппы>>.}
Предназначен для хранения сведений об учебных группах. Реквизиты:

\begin{itemize}
    \item \texttt{Наименование} --- название группы (строка);
    \item \texttt{Курс} --- номер курса (число);
    \item \texttt{Факультет} --- название факультета (строка).
\end{itemize}

\noindent\textbf{Справочник <<Дисциплины>>.}
Предназначен для хранения перечня учебных дисциплин. Реквизиты:

\begin{itemize}
    \item \texttt{Наименование} --- название дисциплины (строка);
    \item \texttt{КодДисциплины} --- код дисциплины (строка);
    \item \texttt{КоличествоЧасов} --- общее количество часов (число).
\end{itemize}

\noindent\textbf{Справочник <<Семестры>>.}
Предназначен для хранения сведений об учебных периодах. Реквизиты:

\begin{itemize}
    \item \texttt{Наименование} --- название учебного периода (строка);
    \item \texttt{УчебныйГод} --- учебный год (строка);
    \item \texttt{НомерСеместра} --- номер семестра (число);
    \item \texttt{ДатаНачала} --- дата начала семестра (дата);
    \item \texttt{ДатаОкончания} --- дата окончания семестра (дата).
\end{itemize}

\noindent\textbf{Справочник <<ПользователиИС>>.}
Предназначен для хранения учётных записей пользователей информационной
системы. Реквизиты:

\begin{itemize}
    \item \texttt{Наименование} --- ФИО пользователя (строка);
    \item \texttt{Логин} --- имя учётной записи (строка);
    \item \texttt{Роль} --- роль пользователя, перечисление
    \texttt{РолиПользователей} (Студент, Преподаватель, Администратор);
    \item \texttt{Студент} --- ссылка на справочник <<Студенты>>;
    \item \texttt{Преподаватель} --- ссылка на справочник <<Преподаватели>>.
\end{itemize}

\noindent\textbf{Документы.}
Для регистрации основных событий учебного процесса в конфигурации создано
три документа. Их структура приведена ниже.

\noindent\textbf{Документ <<ВводОценок>>.}
Предназначен для регистрации оценок, выставляемых преподавателем.
Реквизиты шапки: \texttt{Преподаватель} (ссылка на справочник
<<Преподаватели>>) и \texttt{Семестр} (ссылка на справочник <<Семестры>>).
Табличная часть \texttt{Оценки} содержит реквизиты \texttt{Студент} (ссылка на
<<Студенты>>), \texttt{Дисциплина} (ссылка на <<Дисциплины>>),
\texttt{ДатаОценки} (дата), \texttt{Оценка} (число) и \texttt{Комментарий}
(строка). Показатели документа: количество выставленных оценок и значения
оценок студентов.

\noindent\textbf{Документ <<КорректировкаОценки>>.}
Служит для изменения ранее выставленных оценок. Реквизиты шапки:
\texttt{Преподаватель}, \texttt{Семестр} и \\
\texttt{ПричинаКорректировки} (строка). Табличная часть
\texttt{Корректировки} содержит реквизиты
\texttt{Студент}, \texttt{Дисциплина}, \texttt{СтараяОценка} (число),
\texttt{НоваяОценка} (число) и \texttt{ДатаОценки} (дата). Показатели
документа: старая и новая оценки, а также количество изменённых записей.

\noindent\textbf{Документ <<НазначениеПреподавателя>>.}
Предназначен для закрепления преподавателей за дисциплинами и учебными
группами. Реквизиты шапки: \texttt{Семестр} (ссылка на <<Семестры>>) и
\texttt{ДатаНазначения} (дата). Табличная часть \texttt{Назначения} содержит
реквизиты \texttt{Преподаватель}, \texttt{Дисциплина} и \texttt{УчебнаяГруппа}.
Показатели документа: количество назначений и перечень закреплённых дисциплин
и групп.

Движения, которые документы формируют по регистрам накопления, приведены в
таблице~\ref{tab:movements}.

\begin{table}[H]
    \centering
    \caption{Таблица \thetable\ --- Движения документов по регистрам накопления}
    \label{tab:movements}
    \footnotesize
    \begin{tabularx}{\textwidth}{@{}l >{\raggedright\arraybackslash}X >{\raggedright\arraybackslash}X@{}}
        \toprule
        Документ & Регистр <<ОценкиСтудентов>> & Регистр <<ИтогиУспеваемости>> \\
        \midrule
        ВводОценок &
        Приход: \texttt{Баллы}~=~\texttt{Оценка},
        \texttt{КоличествоОценок}~=~1 &
        Приход: \texttt{СуммаБаллов}~=~\texttt{Оценка},
        \texttt{КоличествоОценок}~=~1 \\
        КорректировкаОценки &
        Расход: \texttt{Баллы}~=~\texttt{СтараяОценка},
        \texttt{КоличествоОценок}~=~1;\newline
        приход: \texttt{Баллы}~=~\texttt{НоваяОценка},
        \texttt{КоличествоОценок}~=~1 &
        Расход: \texttt{СуммаБаллов}~=~\texttt{СтараяОценка},
        \texttt{КоличествоОценок}~=~1;\newline
        приход: \texttt{СуммаБаллов}~=~\texttt{НоваяОценка},
        \texttt{КоличествоОценок}~=~1 \\
        НазначениеПреподавателя &
        Не формирует движений &
        Не формирует движений \\
        \bottomrule
    \end{tabularx}
    \addtocounter{table}{-1}
\end{table}

Документ <<НазначениеПреподавателя>> не формирует движений, поскольку в
конфигурации не предусмотрен отдельный регистр назначений: сведения о
закреплении преподавателей хранятся непосредственно в документе. Это
ограничение можно устранить при развитии системы, добавив регистр сведений
или регистр накопления назначений.

\noindent\textbf{Регистры накопления.}
Для хранения и анализа оценок создано два регистра накопления. Оба регистра
имеют вид \textbf{Обороты}: оценки не образуют складского остатка, поэтому
требуется накапливать обороты и агрегировать показатели за период. Структура
регистров приведена в таблице~\ref{tab:registers}.

\begin{table}[H]
    \centering
    \caption{Таблица \thetable\ --- Регистры накопления конфигурации}
    \label{tab:registers}
    \scriptsize
    \begin{tabularx}{\textwidth}{@{}l l
        >{\hsize=0.9\hsize\raggedright\arraybackslash}X
        >{\hsize=0.9\hsize\raggedright\arraybackslash}X
        >{\hsize=1.2\hsize\raggedright\arraybackslash}X@{}}
        \toprule
        Регистр & Вид & Измерения & Ресурсы & Источники \\
        \midrule
        ОценкиСтудентов & Обороты &
        \texttt{Студент}, \texttt{Дисциплина}, \texttt{Преподаватель},
        \texttt{Семестр} &
        \texttt{Баллы}, \texttt{КоличествоОценок} &
        <<ВводОценок>>,\newline \mbox{<<КорректировкаОценки>>} \\
        ИтогиУспеваемости & Обороты &
        \texttt{Студент}, \texttt{Дисциплина}, \texttt{Семестр} &
        \texttt{СуммаБаллов}, \texttt{КоличествоОценок} &
        <<ВводОценок>>,\newline \mbox{<<КорректировкаОценки>>} \\
        \bottomrule
    \end{tabularx}
    \addtocounter{table}{-1}
\end{table}

\noindent\textbf{Регистр <<ОценкиСтудентов>>.}
Регистр предназначен для хранения информации о выставленных студентам оценках
с возможностью анализа по дисциплинам, преподавателям и семестрам. Измерения
регистра: \texttt{Студент}, \texttt{Дисциплина}, \texttt{Преподаватель},
\texttt{Семестр}. Ресурсы: \texttt{Баллы} (число) и \texttt{КоличествоОценок}
(число). Источники формирования движений: документ <<ВводОценок>> добавляет
новые оценки, документ <<КорректировкаОценки>> фиксирует изменение ранее
зарегистрированных значений.

\noindent\textbf{Регистр <<ИтогиУспеваемости>>.}
Регистр предназначен для хранения агрегированных показателей успеваемости
студентов и быстрого расчёта среднего балла. Измерения регистра:
\texttt{Студент}, \texttt{Дисциплина}, \texttt{Семестр}. Ресурсы:
\texttt{СуммаБаллов} (число) и \texttt{КоличествоОценок} (число). Источники
формирования движений: документ <<ВводОценок>> увеличивает сумму баллов и
количество оценок, документ <<КорректировкаОценки>> изменяет эти показатели с
учётом старого и нового значений. Средний балл рассчитывается как отношение
суммы баллов к количеству оценок:
\[
    \text{Средний балл} = \frac{\texttt{СуммаБаллов}}{\texttt{КоличествоОценок}}.
\]

При корректировке в регистры записываются две записи: расход со старой оценкой
и приход с новой. Для регистра <<ОценкиСтудентов>> обе записи сохраняются в
истории, а для регистра <<ИтогиУспеваемости>> при расчёте итогов учитывается
разница между приходом и расходом, поэтому старая оценка заменяется новой без
двойного счёта.

\noindent\textbf{Иерархическая структура конфигурации.}
В конфигураторе объекты организованы следующим образом:

\begin{itemize}
    \item \textbf{Справочники}:
    \begin{itemize}
        \item Студенты;
        \item Преподаватели;
        \item УчебныеГруппы;
        \item Дисциплины;
        \item Семестры;
        \item ПользователиИС.
    \end{itemize}
    \item \textbf{Документы}:
    \begin{itemize}
        \item ВводОценок;
        \item КорректировкаОценки;
        \item НазначениеПреподавателя.
    \end{itemize}
    \item \textbf{Регистры накопления}:
    \begin{itemize}
        \item ОценкиСтудентов;
        \item ИтогиУспеваемости.
    \end{itemize}
    \item \textbf{Перечисления}:
    \begin{itemize}
        \item РолиПользователей (Студент, Преподаватель, Администратор).
    \end{itemize}
\end{itemize}

Разработанная структура обеспечивает хранение сведений об участниках учебного
процесса, регистрацию оценок, корректировку результатов и расчёт показателей
успеваемости студентов.

\subsection{Этап 2. Реализация функционала автоматизации}

На втором этапе реализована программная логика конфигурации. Код размещён в
модулях форм и модулях объектов документов. Для уменьшения числа ошибок
проверки выполняются как на клиенте при вводе, так и на сервере при проведении
документа.

\noindent\textbf{Экранные формы.}
В форме документа <<ВводОценок>> для табличной части \texttt{Оценки}
обрабатывается событие \texttt{ОценкаПриИзменении}. При вводе значения вне
диапазона от 1 до 5 выводится предупреждение, а оценка сбрасывается в
минимально допустимое значение. Обработчик приведён ниже.

{\setstretch{1.0}\small
\begin{verbatim}
&НаКлиенте
Процедура ОценкиОценкаПриИзменении(Элемент)

    Если Элемент.ТекущиеДанные.Оценка < 1
        Или Элемент.ТекущиеДанные.Оценка > 5 Тогда

        Предупреждение("Оценка должна быть в диапазоне от 1 до 5.");
        Элемент.ТекущиеДанные.Оценка = 1;

    КонецЕсли;

КонецПроцедуры
\end{verbatim}
}

Автоподстановка данных реализована в формах документов: при добавлении строки
табличной части дата оценки заполняется текущей датой, преподаватель и семестр
берутся из шапки документа, а в документе <<КорректировкаОценки>> старая оценка
подставляется по выбранным студенту и дисциплине. Автоматический расчёт
показателей выполняется при проведении: ресурсы \texttt{СуммаБаллов} и
\texttt{КоличествоОценок} формируются программно, а средний балл рассчитывается
как их отношение.

\noindent\textbf{Механизмы проведения документов.}
Проведение документов реализовано в модуле объекта документа в процедуре
\texttt{ОбработкаПроведения}. Для документа <<ВводОценок>> выполняется обход
табличной части \texttt{Оценки}; для каждой строки создаются движения в обоих
регистрах накопления. Фрагмент обработчика приведён ниже.

{\setstretch{1.0}\small
\begin{verbatim}
Процедура ОбработкаПроведения(Отказ, РежимПроведения)

    Движения.ОценкиСтудентов.Записывать = Истина;
    Движения.ИтогиУспеваемости.Записывать = Истина;

    Для Каждого СтрокаОценки Из Оценки Цикл

        Движение = Движения.ОценкиСтудентов.Добавить();
        Движение.ВидДвижения = ВидДвиженияНакопления.Приход;
        Движение.Период = Дата;
        Движение.Студент = СтрокаОценки.Студент;
        Движение.Дисциплина = СтрокаОценки.Дисциплина;
        Движение.Преподаватель = Преподаватель;
        Движение.Семестр = Семестр;
        Движение.Баллы = СтрокаОценки.Оценка;
        Движение.КоличествоОценок = 1;

        Движение = Движения.ИтогиУспеваемости.Добавить();
        Движение.ВидДвижения = ВидДвиженияНакопления.Приход;
        Движение.Период = Дата;
        Движение.Студент = СтрокаОценки.Студент;
        Движение.Дисциплина = СтрокаОценки.Дисциплина;
        Движение.Семестр = Семестр;
        Движение.СуммаБаллов = СтрокаОценки.Оценка;
        Движение.КоличествоОценок = 1;

    КонецЦикла;

КонецПроцедуры
\end{verbatim}
}

В документе <<КорректировкаОценки>> для каждой строки табличной части сначала
формируется расходная запись со старой оценкой, затем приходная запись с новой
оценкой. Аналогичные пары записей создаются в регистре
<<ИтогиУспеваемости>>: расход уменьшает сумму баллов на старую оценку, приход
увеличивает её на новую. В результате история изменений сохраняется, а итоговые
показатели остаются корректными.

\noindent\textbf{Контроль и бизнес-логика.}
Перед формированием движений в процедуре \texttt{ОбработкаПроведения}
выполняются проверки, предотвращающие проведение некорректного документа.
Проверяется заполнение преподавателя и семестра, наличие хотя бы одной строки
в табличной части, заполнение студента и дисциплины в каждой строке, а также
допустимый диапазон оценки. В предметной области учёта успеваемости эти
проверки являются аналогом контроля достаточности товара на складе: они не
позволяют провести документ с неполными или недостоверными данными. Фрагмент
проверок приведён ниже.

{\setstretch{1.0}\small
\begin{verbatim}
Если Не ЗначениеЗаполнено(Преподаватель) Тогда
    Сообщить("Не указан преподаватель.");
    Отказ = Истина;
    Возврат;
КонецЕсли;

Если Не ЗначениеЗаполнено(Семестр) Тогда
    Сообщить("Не указан семестр.");
    Отказ = Истина;
    Возврат;
КонецЕсли;

Если Оценки.Количество() = 0 Тогда
    Сообщить("Не добавлено ни одной оценки.");
    Отказ = Истина;
    Возврат;
КонецЕсли;

// Проверки внутри цикла по строкам табличной части
Если Не ЗначениеЗаполнено(СтрокаОценки.Студент)
    Или Не ЗначениеЗаполнено(СтрокаОценки.Дисциплина) Тогда

    Сообщить("Не заполнен студент или дисциплина.");
    Отказ = Истина;
    Возврат;

КонецЕсли;

Если СтрокаОценки.Оценка < 1
    Или СтрокаОценки.Оценка > 5 Тогда

    Сообщить("Оценка должна быть в диапазоне от 1 до 5.");
    Отказ = Истина;
    Возврат;

КонецЕсли;
\end{verbatim}
}

При обнаружении ошибки устанавливается признак \texttt{Отказ}~=~\texttt{Истина},
выводится сообщение и выполняется возврат из процедуры. Документ не проводится,
движения по регистрам не записываются. Проверка на клиенте предупреждает
пользователя сразу при вводе, а проверка при проведении защищает от обхода
формы и гарантирует целостность данных.

\noindent\textbf{Итоги реализации.}
Соответствие выполненной работы требованиям задания приведено в
таблице~\ref{tab:requirements}.

\begin{table}[H]
    \centering
    \caption{Таблица \thetable\ --- Соответствие требований задания и реализации}
    \label{tab:requirements}
    \footnotesize
    \begin{tabularx}{\textwidth}{@{}>{\raggedright\arraybackslash}X l >{\raggedright\arraybackslash}X@{}}
        \toprule
        Требование & Результат & Реализация \\
        \midrule
        Не менее 5 справочников & 6 & Студенты, Преподаватели,
        УчебныеГруппы, Дисциплины, Семестры, ПользователиИС \\
        Не менее 3 документов & 3 & ВводОценок, КорректировкаОценки,
        НазначениеПреподавателя \\
        Не менее 2 регистров накопления & 2 & ОценкиСтудентов,
        ИтогиУспеваемости (вид <<Обороты>>) \\
        Автоматический расчёт & Выполнен & Ресурсы \texttt{СуммаБаллов} и
        \texttt{КоличествоОценок} рассчитываются при проведении; средний балл
        вычисляется как отношение ресурсов \\
        Автоподстановка данных & Выполнена & Дата, преподаватель и семестр
        подставляются из шапки; старая оценка --- по студенту и дисциплине \\
        Механизмы проведения & Реализованы & Движения формируются в
        \texttt{ОбработкаПроведения} документов <<ВводОценок>> и
        <<КорректировкаОценки>> \\
        Контроль и бизнес-логика & Реализованы & Проверки заполнения
        реквизитов, наличия строк и диапазона оценки 1--5 \\
        \bottomrule
    \end{tabularx}
    \addtocounter{table}{-1}
\end{table}

\section{Вывод}

В ходе работы разработана конфигурация «1С:Предприятие 8.3» для автоматизации
учёта успеваемости студентов. На этапе проектирования созданы шесть
справочников (<<Студенты>>, <<Преподаватели>>, <<УчебныеГруппы>>,
<<Дисциплины>>, <<Семестры>>, <<ПользователиИС>>), три документа
(<<ВводОценок>>, <<КорректировкаОценки>>, <<НазначениеПреподавателя>>) и два
регистра накопления (<<ОценкиСтудентов>> и <<ИтогиУспеваемости>>), а также
перечисление <<РолиПользователей>>. Структура конфигурации отражает предметную
область и превышает минимальные требования задания по числу справочников.

На этапе программирования реализованы проверка оценки в форме документа,
автоподстановка данных, автоматический расчёт ресурсов регистров и механизмы
проведения документов. Документ <<ВводОценок>> формирует приходные движения в
обоих регистрах, документ <<КорректировкаОценки>> фиксирует расход по старой
оценке и приход по новой, а документ <<НазначениеПреподавателя>> сохраняет
сведения о закреплении преподавателей без движений по регистрам. Контроль
заполнения реквизитов, наличия строк и диапазона оценки 1--5 предотвращает
проведение некорректных документов и защищает данные от ошибок пользователя.

Таким образом, цель работы достигнута: конфигурация обеспечивает хранение
нормативно-справочной информации, регистрацию и корректировку оценок, расчёт
показателей успеваемости и работу без логических ошибок. В качестве направления
развития можно добавить регистр назначений преподавателей, отчёты об
успеваемости и разграничение прав доступа пользователей.

\end{document}
